\section{Введение: как устроен курс}

\subsection*{О чём эта книга}

Эта книга учит программировать квадрокоптеры «Геоскан Пионер». Их два, и программируются они по-разному:

\begin{center}
\small
\begin{tabularx}{\linewidth}{@{}>{\bfseries}l X X@{}}
\toprule
 & Пионер Базовый & Пионер Мини~2 \\
\midrule
Языки & Lua (на борту), Python (на компьютере) & Python (на борту), блоки Pioneer Bricks \\
Библиотеки & \texttt{pioneer\_sdk} & \texttt{pioneer\_sdk2}, \texttt{pioneer\_rknn}, \texttt{rclpy} (ROS~2), \texttt{pioneer\_swarm} \\
Нейросети & нет & нейропроцессор на борту \\
Рой дронов & с компьютера (глава~\ref{sec:python-swarm}) & Swarm Kit на борту (часть V) \\
Где в книге & части I--III & части IV и V \\
\bottomrule
\end{tabularx}
\end{center}

Базовый дрон --- основа курса: на нём от первой строчки на Lua идёт путь до автономной миссии, в которой дрон сам взлетает, облетает точки, реагирует на нештатные ситуации и садится. Та же программа на Python выходит за пределы симулятора --- к камере, маркерам и групповым полётам. Пионер Мини~2 --- следующая ступень: программа работает на самом дроне, а рядом есть нейропроцессор. Для него же написаны две заключительные части курса: ROS~2 (пакет \texttt{gs\_core}), в котором дрон устроен как набор узлов, и Swarm Kit --- рой, где дроны находят друг друга и работают вместе без центра. Эти главы лучше читать после основ Python.

Главный принцип курса --- \textbf{каждый пример запускается}. Программы для Lua и \texttt{pioneer\_sdk} лежат в папке \texttt{examples/} рядом с текстом и автоматически проверяются в симуляторе: взлетают, пролетают свои точки и садятся. Если такой пример не совпадает с тем, что вы видите на экране, --- это ошибка книги, а не ваша. Примеры Мини~2 (\texttt{examples/python-onboard/}) написаны для борта дрона: их проверяют на имитации библиотек по следу команд (примеры ROS~2 и роя --- на эмуляции \texttt{rclpy} и \texttt{pioneer\_swarm} в симуляторе), а какие из них идут и в симуляторе, сказано в самих главах. На настоящем дроне полётные программы Мини~2, ROS~2 и роя авторы не запускали.

\subsection*{Где писать код}

Для занятий нужен только браузер: веб-симулятор \url{https://simulator.sakoryagin.ru} показывает полёт в 3D, умеет Lua, Python и блоки, содержит обе модели дрона и использует те же API, что настоящие Пионеры. Настоящий дрон понадобится ближе к концу курса --- для камеры, компьютерного зрения и бортовых программ. Как запускать программы в симуляторе и на дронах, объясняет глава \ref{sec:start}.

\subsection*{Как читать главы}

Каждая глава построена одинаково:
\begin{enumerate}
    \item \textbf{Цели} --- что вы будете уметь после главы.
    \item \textbf{Миссии симулятора} --- зелёная врезка: какие миссии режима \ui{Обучение} закрепляют тему главы (раздел~\ref{sec:learning-mode}).
    \item \textbf{Уроки} --- короткая теория с примерами: программы, которые нужно запустить и сверить результат с комментариями в коде.
    \item \textbf{Задание} --- рамка после каждого урока: что сделать и что должно получиться. Проверять результат можно прямо в симуляторе.
    \item \textbf{Проверьте себя} --- вопросы с ответами.
    \item \textbf{Итоги} --- что запомнить.
\end{enumerate}

Рисунок~\ref{fig:course-loop} показывает этот цикл по шагам.

\begin{tikzfig}
\begin{adjustbox}{max width=\linewidth}\begin{tikzpicture}[fig]
\node[fblue, text width=28mm, minimum height=19mm] (a) at (0,0) {\textbf{1. Урок} \\ \footnotesize короткая теория и пример программы};
\node[forange, text width=28mm, minimum height=19mm] (b) at (4.7,0) {\textbf{2. Запуск} \\ \footnotesize пример в симуляторе, сверка с комментариями};
\node[forange, text width=28mm, minimum height=19mm] (c) at (9.4,0) {\textbf{3. Задание} \\ \footnotesize рамка после урока и ожидаемый результат};
\node[fgreen, text width=28mm, minimum height=19mm] (d) at (14.1,0) {\textbf{4. Миссия} \\ \footnotesize проверка полёта в режиме «Обучение»};
\draw[farrow] (a) -- (b);
\draw[farrow] (b) -- (c);
\draw[farrow] (c) -- (d);
\draw[fcmd] (d.south) -- ++(0,-1.0) -| (a.south);
\node[flab] at (7.05,-1.95) {следующий урок};
\end{tikzpicture}\end{adjustbox}
\captionof{figure}{Рабочий цикл по главе: прочитать урок, запустить пример, выполнить задание, пройти миссию симулятора и перейти к следующему уроку. Главу закрывают «Проверьте себя» и «Итоги».}\label{fig:course-loop}
\end{tikzfig}


В тексте встречаются рамки:
\begin{itemize}
    \item \textcolor{GeoskanBlue}{\faInfoCircle}~\textbf{Информация} --- пояснение или ссылка на другой раздел;
    \item \textcolor{red!70!black}{\faExclamationTriangle}~\textbf{Внимание} --- типичная ошибка или вопрос безопасности;
    \item \textcolor{GeoskanOrange}{\faLightbulb}~\textbf{Подсказка} --- приём, который экономит время;
    \item \textcolor{GeoskanDarkBlue}{\faGavel}~\textbf{Правило} --- соглашение, которого курс придерживается всегда;
    \item \textcolor{GeoskanOrange}{\faPencilAlt}~\textbf{Задание} --- работа для самостоятельного выполнения;
    \item \textcolor{TaskTeal}{\faGamepad}~\textbf{Миссии симулятора} --- связь главы с режимом \ui{Обучение}.
\end{itemize}

\subsection*{Дорожная карта}

\begin{enumerate}
    \item \textbf{Подготовка.} Интерфейс симулятора, режим \ui{Обучение}, первая программа, настоящие дроны и техника безопасности.
    \item \textbf{Часть I. Lua} (Пионер Базовый). Язык (переменные, условия, циклы, функции) $\to$ светодиоды и таймеры $\to$ события автопилота и автономные миссии.
    \item \textbf{Часть II. Python.} Основы языка $\to$ \texttt{pioneer\_sdk} (Базовый): подключение, светодиоды, взлёт, полёт по точкам $\to$ камера, маркеры ArUco, трекинг, рой.
    \item \textbf{Часть III. Профессиональная навигация.} Теория визуальных маркеров, калибровка камеры, карты, фильтрация и точная посадка (на \texttt{pioneer\_sdk}).
    \item \textbf{Часть IV. Пионер Мини~2.} \texttt{pioneer\_sdk2} (глава~\ref{sec:sdk2}), нейросети \texttt{pioneer\_rknn} (глава~\ref{sec:rknn}) и блоки Pioneer Bricks (глава~\ref{sec:bricks}).
    \item \textbf{Часть V. ROS~2 и рой на Пионере.} Узлы, топики и сервисы \texttt{gs\_core} (глава~\ref{sec:ros2}): полётный цикл, системы координат \texttt{map} и \texttt{base\_link}, телеметрия, \texttt{cmd\_vel}, параметры, нагрузка, миссия-узел $\to$ Swarm Kit (глава~\ref{sec:swarm}): соседи, сообщения, телеметрия соседей, строй, групповой проект.
    \item \textbf{Приложения.} Справочник Lua API, шпаргалка по языкам и библиотекам (включая ROS~2 и рой), настройка пульта, протокол безопасности.
\end{enumerate}

Lua и Python в курсе идут в одном порядке: сначала язык, потом светодиоды, потом полёт. Поэтому часть про Python можно читать и отдельно, но полезнее --- после Lua: событийная модель автопилота, разобранная на Lua, объясняет, почему Python-программа «ждёт точку» именно так. Главы про Мини~2 опираются на Python-основы и на полёт по точкам из \texttt{pioneer\_sdk}; блоки Bricks можно начать и раньше, не дожидаясь текстового Python. Часть V строится на главе~\ref{sec:sdk2}: полёт \texttt{pioneer\_sdk2} объясняет, что делают сервисы \texttt{gs\_core}, а рой дополняет их общением между дронами.

\subsection*{Планируемые результаты}

После курса вы сможете:
\begin{itemize}
    \item писать читаемый код на Lua и Python и объяснять, что делает каждая строка;
    \item управлять бортовой периферией --- светодиодами и таймерами --- без блокировки программы;
    \item строить программу полёта как конечный автомат: команда $\to$ событие $\to$ следующая команда;
    \item предусматривать аварийные ситуации: удар, разряд батареи, потеря связи;
    \item переносить программу из симулятора на настоящий дрон, учитывая отличия моделей (Lua и \texttt{pioneer\_sdk} у Базового, \texttt{pioneer\_sdk2} у Мини~2);
    \item подключить к полёту компьютерное зрение и нейросети: маркеры ArUco, распознавание людей и текста;
    \item написать узел ROS~2 для Пионера (сервисы, топики, параметры) и группу дронов, которая общается сообщениями и следует за ведущим в строю.
\end{itemize}
